% !PRINTED 492
\weberpage{492}
%<<CONT
not identical with $\psi$. One can therefore (by \S.\,143, 1.) give $t$ such a rational value that $\psi$ is different from all functions formed after the manner of $\psi_1$.

It is self-evident that these considerations also hold when independent variables take the place of the $\alpha$.

A function which admits all the permutations of a group $P$, while it changes under all permutations not belonging to $P$, is called a function \gesperrt{belonging to the group $P$}.

To the group $P_0$, which comprises \gesperrt{all} permutations of $m$ figures, belong the symmetric functions. This group is therefore also called the \gesperrt{symmetric group}.

A function $\varrho$ which belongs to the group consisting of the single identical permutation can serve as an unambiguous way of denoting the permutations; for there is only one permutation $\pi'$ by which $\varrho$ passes into $\varrho'$. We can therefore use for this permutation the symbol
\[
\pi' = (\varrho, \varrho')
\]
so that by
\[
(\varrho, \varrho), \quad (\varrho, \varrho'), \quad (\varrho, \varrho'') \dots
\]
all permutations of the $m$ figures are unambiguously denoted.

\begin{center}
\S.\,153.

\gesperrt{Decomposition of Permutations into Transpositions\\
and into Cycles.}
\end{center}

We have already mentioned, in the second section, on the occasion of the determinants, the composition of arbitrary permutations from a sequence of transpositions, that is, the permutation of only two figures. Such transpositions we denote by placing the figures in question side by side, e.g.\ by $(0,1)$ the permutation of the figures 0 and 1.

A transposition, repeated twice, leads to the identity, so that every transposition is opposite to itself. Now if
\[
\pi = \begin{pmatrix} 0, & 1, & 2 \dots m-1 \\ a_0, & a_1, & a_2 \dots a_{m-1} \end{pmatrix}
\]
is an arbitrary permutation of $m$ figures, then
\[
\pi(m-1, a_{m-1}) = \pi_1
\]
is a permutation which leaves $m-1$ unchanged, and which therefore by
%CONT>>
